% 第 51 章　原子模型的兴衰
\chapter{原子模型的兴衰}

前面几章里，量子力学已经亮过相：能量可以一份一份地来，光可以表现得像粒子，粒子也可以表现得像波。但历史上，这些观念不是从黑板上长出来的，而是被同一个东西逼出来的——原子。二十世纪初的物理学家手里握着一堆关于原子的实验事实，却发现任何一个“常识模型”都解释不了它们。这一章我们就回到那个时代，看三种原子模型怎样轮番登场、轮番倒下，又怎样在废墟上留下“能级”这个直到今天还在用的概念。读完你会发现：光谱学是这场革命的主战场，而你家厨房里就能重现第一声枪响。

\step{反常识开场}

\begin{mainline}
炒菜的时候如果不小心把盐撒进燃气灶的蓝色火苗里，你会看到火苗“腾”地一下变成明亮的黄色。不是橙黄，不是金黄，是一种很纯的黄，像有人往火焰里挤了一滴黄色颜料。烟花师傅靠的就是这套把戏：硝酸锶让烟花变红，钡盐让它变绿，铜盐让它变蓝。夜晚的街道上，老式路灯发出熟悉的橙黄光，霓虹灯管发出橙红光——它们的颜色都不是染料染出来的，而是气体自己在“报名字”。

现在来做一次对比。白炽灯的光透过棱镜，展开成一条平滑的彩虹，从红到紫一个颜色都不缺；太阳光也是一样。可如果把氢气封进玻璃管，加上高电压让它发光，再透过棱镜去看，你看到的不是彩虹，而是漆黑背景上几条孤零零的亮线：一条鲜红，一条青蓝，一两条暗紫，其余位置什么都没有。发光的气体挑剔得不可思议：它只肯发出某几种颜色，仿佛原子手里拿着一张严格的颜色清单。

更蹊跷的事情发生在太阳身上。1814 年，夫琅禾费用当时最好的棱镜细看太阳光谱，发现彩虹里嵌着几百条细细的暗线，像有人用针在彩带上扎出一排排小孔。后来人们把氢放电管的亮线位置拿去和太阳光谱的暗线位置比对，结果发现：氢管里发亮的那几个波长，在太阳那里恰好就是缺掉的那几个波长，位置分毫不差。同一种元素，在一种场合“出席”，在另一种场合“缺席”，而且出席和缺席的名单完全一致。这哪里是光的颜色，这分明是原子留下的指纹。问题是：指纹是谁按上去的？原子内部到底藏着什么机关，能让它对颜色如此挑剔？
\end{mainline}

\step{先猜一猜}

\begin{mainline}
在揭开机关之前，先把你自己的直觉摆上桌面。下面四个说法，凭第一感觉判断对错，最好写下来。

第一，不管什么东西，烧得足够热都会发出差不多的光，颜色只由温度决定，跟材料关系不大。

第二，原子就是实心小球，不可分割；化学反应无非是小球们重新排队，小球本身永远不变。

第三，假如原子内部真的有一个带正电的核和一个带负电的电子，那么电子被核吸引，理应像地球绕太阳那样，沿着一条稳定的轨道一直转下去——这是天体物理交给我们的常识。

第四，各种元素发出的谱线位置杂乱无章，是历史的偶然；想从一堆波长数字里找出整齐的规律，多半是白费力气。

读完本章再回头对照。第一条对了一半：它适用于烧热的固体和稠密物质，却在稀薄气体面前输得很惨。第二条被 1897 年的一项发现直接推翻。第三条错得最要命——正是这条“天体常识”把二十世纪初的物理学逼进了死胡同，本章一半的戏都在它身上。第四条错得最冤枉：谱线背后的规律不但存在，而且整齐到初中生用四则运算就能验证。物理学家当年也不敢相信，一位中学数学老师先把答案猜中了。
\end{mainline}

\step{动手实验}

\begin{experiment}[用一张废光盘给光“分户口”]
你不需要棱镜，也不需要实验室。找一张废旧 CD 或 DVD，它表面那圈彩虹本身就出卖了秘密：盘片上刻着间距约 \(1.6\) 微米的螺旋轨道，这个间距只有头发丝直径的几十分之一，刚好是可见光波长的两三倍。一排等间距的细缝或细槽会把不同颜色的光弯向不同的角度，这种装置叫光栅——你的废光盘就是一块现成的反射光栅。

关灯，只留光源。把光盘的反光面朝向光源，调整角度，让彩色光带落在你眼睛或手机镜头里。依次观察：

蜡烛火焰——一条平滑的彩虹，红橙黄绿蓝靛紫连在一起，没有缺口。

白炽灯泡——同样平滑的彩虹。

日光灯管或节能灯——底色是连续的彩虹，但上面压着几条格外亮的细线，特别是绿色和蓝紫色位置。那是灯管里汞蒸气的“指纹”透过荧光粉露了出来。

手机白屏幕——三块宽宽的色带，红、绿、蓝各一摊，中间还有暗区。你的“白色”其实是三种颜色拼出来的错觉，光盘一眼就看穿了它。

最后一步最精彩。在蜡烛火焰上方撒几粒食盐（或者滴几滴浓盐水），火焰变黄。用光盘去看这束黄光：彩虹的其他位置几乎消失，只剩黄色位置上一根格外明亮的细线。你刚刚完成了一次货真价实的光谱分析：从光里认出了钠原子。天体物理学家对星光做的事，和这件事在原理上没有区别。
\end{experiment}

\begin{mainline}
这个实验顺手推翻了一个朴素的假设：发光体的光谱是“连续底色加细节”。事实上恰好反过来——每种发光体都有自己的“户籍”，连续光谱和线状光谱是两种根本不同的发光方式。光盘把光按波长排开，等于把原子交上来的账本一页页摊开给你看。接下来的问题是：账本是谁记的？
\end{mainline}

\step{旧模型遇到麻烦}

\begin{mainline}
麻烦要一层一层揭开，因为历史就是一层一层撞墙的。

第一堵墙之前，是化学家的天下。1808 年，道尔顿把古希腊人的猜想变成定量科学：化学反应里各物质的质量比是固定的简单整数比，最自然的解释就是物质由不可分割的小球组成，反应只是小球重新组合。定比定律、倍比定律都被这套“实心小球”模型解释得漂漂亮亮。它是好模型——直到电来了。

1897 年，汤姆孙研究真空管里的阴极射线。这种从负极射出的神秘“光线”能被电场和磁场掰弯，弯曲方向说明它带负电，弯曲程度能算出一个关键数字：粒子的电荷与质量之比。决定性的发现是——换掉电极的金属材料，换掉管子里残存的气体，这个比值纹丝不动，永远是氢离子的约一千八百倍。也就是说，所有物质里都藏着同一种带负电的小颗粒，它比最轻的原子还轻近两千倍。电子就这样出场了。实心小球模型当场作废：原子可分，而且每个原子肚子里都有电子。

可是原子整体是电中性的，电子带负电，正电在哪里？汤姆孙给出了第二代的答案：原子像一块枣糕，正电荷是均匀分布的“糕”，电子是嵌在里面的“枣”。它像什么？它像果冻里冻着的果粒，整个结构稳定、电中性，果粒还能在平衡位置附近微微振动——振动恰好能发光，定性上似乎连光谱都有着落。它不像什么？它不像一个“空”原子——在枣糕模型里，原子内部被正电荷填得满满当当，任何打进来的带电粒子都应该受到四面八方大致抵消的弱电场，顶多偏转一点点。

1909 年，盖革和马斯登在卢瑟福的实验室里检验这个预言。他们用镭放出的 \(\alpha\) 粒子（带两份正电、比电子重七千多倍的“炮弹”）轰击只有几百层原子厚的金箔，然后在四周观察粒子飞向哪里。结果绝大多数 \(\alpha\) 粒子径直穿过，这倒不意外；真正吓人的是例外——大约每八千个粒子里就有一个偏转超过九十度，有的甚至几乎原路弹了回来。按枣糕模型估算，靠无数次小偏转累积出一次大角度弹回，概率小到小数点后几千个零，等到宇宙热寂也等不来一次。卢瑟福后来说：“这就像你朝一张卫生纸发射一枚十五英寸的炮弹，炮弹却弹回来打中你自己。”直觉在这里输得精光：原子内部不是满满当当，而是空空荡荡。正电荷和几乎全部质量被压缩在一个直径不到原子万分之一的核心里。1911 年，卢瑟福提出核式模型：中心是原子核，电子在外围。图里画的就是这场实验的示意。
\end{mainline}

\begin{figure}[htbp]
\centering
\begin{tikzpicture}[scale=1.05]
  \fill[red!70!black] (0,0) circle (2.2pt);
  \node[below=3pt] at (0,-0.05) {\scriptsize 金核};
  \draw[thick,-{Stealth}] (-4.2,1.7) .. controls (-1,1.62) and (1,1.55) .. (3.2,1.95);
  \draw[thick,-{Stealth}] (-4.2,0.55) .. controls (-0.9,0.5) and (0.4,0.75) .. (2.6,2.5);
  \draw[thick,-{Stealth}] (-4.2,0.08) .. controls (-0.45,0.08) and (-0.35,0.3) .. (-2.1,1.35);
  \node[left] at (-4.2,0.8) {\scriptsize \(\alpha\) 粒子};
\end{tikzpicture}
\caption{卢瑟福散射示意：离核远的 \(\alpha\) 粒子几乎直行，离得近的偏转很大，正对核冲去的会被弹回。绝大多数粒子穿过说明原子内部几乎是空的。}
\end{figure}

\begin{mainline}
先替核式模型算一笔真实的账。金箔实验里弹回比例约为八千分之一，结合金箔大约一千层原子的厚度，可以反推出核的横截面积只占原子横截面积的十亿分之一量级：原子直径约 \(10^{-10}\) 米，原子核直径只有约 \(10^{-14}\) 米。把原子放大到一座足球场，原子核不过是场地中央的一粒玻璃弹珠，却集中了原子百分之九十九点九七以上的质量。你以为墙是实心的，墙其实几乎是真空——手感上的“实心”来自电磁排斥，而不是物质填满。

然而核式模型刚成立就撞上两堵更硬的墙，而且这两堵墙是旧物理学自己的承重墙。

第一堵墙：原子凭什么活着？电子带负电，核带正电，异号电荷相互吸引。电子想不掉进核里，唯一的经典出路是绕核高速运动，让运动的效果抵住吸引——就像地球绕太阳。可是圆周运动是加速运动，而麦克斯韦电磁理论有一条铁推论：加速运动的带电粒子必然辐射电磁波，把能量带走。算算看：绕核的电子一边转一边发光，能量不断流失，轨道一圈圈收缩，最后螺旋坠毁在核上。整个过程要多久？经典电磁学给出的答案是约 \(10^{-11}\) 秒，一百亿分之一秒量级。按照这个理论，原子诞生的瞬间就该死掉，而你身体里的氢原子从大爆炸一直活到今天。

第二堵墙：就算电子暂时转着，光谱也对不上。电子螺旋下坠的过程中，转动的频率连续地变化，辐射的颜色就该连续地扫过整条彩虹——可实验里是几条孤零零的线。核式模型连开场那撮盐的黄光都解释不了。

到 1911 年，局面变得极其难堪：实验逼着我们相信核式结构，理论却断言这种结构活不过一百亿分之一秒。物理学不是缺一个更好的答案，而是所有旧答案都被判了死刑。该发明新东西了。
\end{mainline}

\step{发明新概念}

\begin{mainline}
1913 年，二十八岁的玻尔在卢瑟福的实验室里做了那件物理学家在绝境中才会做的事：既然旧规则集体失效，那就立几条新规则，只要它们能同时安抚实验和理论。注意，下面三条不是从旧理论“推导”出来的，是发明的；发明的价值要用它算出来的数字来偿还。

第一条，定态假设。原子存在一组特殊的状态，称为定态。电子处在定态中时，不辐射能量。这条规则明目张胆地顶撞麦克斯韦——玻尔自己也知道。他的态度是：麦克斯韦理论在宏观世界战功赫赫，但谁也没验证过它在原子内部依然成立；那就先记账，看结果。

第二条，跃迁假设。电子不连续地“滑”过能量，只能在定态之间整体地跳跃。从高能级跳到低能级，原子放出一个光子；从低能级跳到高能级，必须先吸收一个光子。光子的能量严格等于两级之差：
\[
h\nu = E_{\text{高}}-E_{\text{低}}
\]
其中 \(h\) 是普朗克常量，\(\nu\) 是光的频率。这条规则的杀伤力在于：既然能级是特定的一些值，能量差就是特定的一些值，发出的光就只有特定的几种颜色。分立谱线不再是谜，而是规则的直接后果。原子对颜色的“挑剔”，其实是它对能量的挑剔。

第三条，量子化条件。哪些定态被允许？玻尔给出的判据是：电子的角动量只能是约化普朗克常量 \(\hbar=h/(2\pi)\) 的整数倍，
\[
L = n\hbar,\qquad n=1,2,3,\dots
\]
为什么是角动量？因为要挑出“哪些轨道合法”，需要一个带整数的条件；而普朗克常量的量纲恰好是角动量的量纲，整数倍是最简单的选择。这个理由当年听起来像凑数。十年后德布罗意回头看了一眼，才发现它并不任性：如果电子本身是一种波，那么绕核一周，波必须首尾相接、自我叠加而不是自我抵消，这就要求轨道周长恰好是波长的整数倍——把德布罗意波长代进去，出来的正是 \(L=n\hbar\)。玻尔凭直觉立下的规矩，恰好是波在圆环上的驻波条件。

现在履行本书的规矩，声明这个模型像什么、不像什么。它像一架梯子：电子只能站在梯档上，不能悬在半空；跳下一档就放出一笔固定的能量。它不像真实的楼梯：电子不是一个站在某一级的小球，“在第几级”是整个量子态的标签，不是空间位置。它也不像被行政命令规定的行星轨道：定态背后有波的本性撑腰，而这条腰要到下一章薛定谔方程出场才真正露出来。在此之前，请把它当作一张“先记账、后还钱”的欠条。
\end{mainline}

\step{一句话模型}

\begin{oneline}
原子的能量是一架梯子，不是一段斜坡；电子只能在梯档之间跳，每跳一次就按档距兑换一个光子——所以光谱是原子的指纹，档距是原子的身份证号。
\end{oneline}

\begin{mainline}
用这句话重述前面所有现象：焰色反应的纯黄，是钠原子梯子上某两档之间的固定档距；氢放电管的几条亮线，是氢梯子上少数几档之间的距离；太阳光谱的暗线，是太阳大气里的原子按自己的档距把路过的光“吃”掉了。同一张梯子，往下跳是亮线，往上跳是暗线，波长分毫不差——因为档距是同一笔账。
\end{mainline}

\step{数学翻译器}

\begin{mainline}
现在把直觉翻译成数字。先建经典账本。电子绕核做圆周运动时，提供向心效果的是核对电子的库仑引力——向心不是一种新力，是库仑力扮演的角色。库仑力的大小是
\[
F=\frac{ke^{2}}{r^{2}}
\]
它与万有引力一样是平方反比，所以账本结构和行星轨道几乎一样。圆周运动要求 \(mv^{2}/r=ke^{2}/r^{2}\)，配上电势能 \(-ke^{2}/r\)，总能量为
\[
E=\frac{1}{2}mv^{2}-\frac{ke^{2}}{r}=-\frac{ke^{2}}{2r}
\]
请盯着这个负号看一会儿。总能量是负的，意思是电子被“绑”在核旁：要把它彻底拆走，必须额外补给它 \(ke^{2}/(2r)\) 的能量。再做两个特殊检查。\(r\) 减小时，总能量更低（更负），可动能反而变大——束缚得越紧，跑得越快，这和卫星降低轨道反而加速是同一笔账，初看反直觉，算完就服气。\(r\to\infty\) 时，\(E\to 0\)，对应电子刚刚挣脱、恰好静止在无穷远处：这就是电离的门槛。

下一步，把玻尔的第三条规则放进来：\(L=mvr=n\hbar\)。整数 \(n\) 一进方程，连续的 \(r\) 立刻被筛成一串离散的值（推导见数学桥层）：
\[
r_{n}=n^{2}a_{0},\qquad a_{0}\approx 5.3\times10^{-11}\ \mathrm{m}
\]
\[
E_{n}=-\frac{13.6\ \mathrm{eV}}{n^{2}}
\]
\(a_{0}\) 叫玻尔半径，\(13.6\) 电子伏特是氢的电离能。检查时间。\(n=1\)：半径 \(0.053\) 纳米，和实验测出的氢原子尺寸吻合；能量 \(-13.6\) 电子伏特，实验上电离氢气需要的电压正是约 \(13.6\) 伏。\(n\to\infty\)：能量升到 \(0\)，半径趋向无穷——梯子的顶端就是电离，级数越大梯子越挤。方向反转检查：从 \(n=2\) 跌到 \(n=1\)，放出 \(10.2\) 电子伏特，对应波长 \(122\) 纳米的紫外光；反过来从 \(n=1\) 升到 \(n=2\)，吸收的还是 \(122\) 纳米。同一档距，亮线暗线同一位置——开场悬念在公式里兑现了。

最漂亮的检验来自可见光。从 \(n=3\) 跌到 \(n=2\)：能量差 \(13.6\times(1/4-1/9)\approx1.89\) 电子伏特，对应波长 \(656\) 纳米——正是氢放电管里那条最亮的红线，光谱学里叫 H\(\alpha\)。再算 \(n\to 2\) 的整族跃迁，发现它们排成的队形恰好就是 1885 年中学数学老师巴耳末从数据里硬凑出来的公式
\[
\frac{1}{\lambda}=R\left(\frac{1}{2^{2}}-\frac{1}{n^{2}}\right),\qquad n=3,4,5,\dots
\]
里德伯常量 \(R\approx1.097\times10^{7}\ \mathrm{m^{-1}}\)。特殊检查：\(n=3\) 给出 \(656\) 纳米；\(n\to\infty\) 时 \(1/\lambda\) 趋向 \(R/4\)，即波长 \(364.6\) 纳米——谱线越挤越密，却永远不会越过这条“系列极限”线，像一队人挤向一堵墙。巴耳末当年是猜中了一个整数游戏，玻尔给出了这个游戏的来历：那两个平方，就是两个能级的编号。
\end{mainline}

\begin{mathbridge}
动手把 \(r_{n}\) 和 \(E_{n}\) 推出来，只需要两条式子联立。库仑力提供向心效果：
\[
\frac{mv^{2}}{r}=\frac{ke^{2}}{r^{2}}\quad\Longrightarrow\quad mv^{2}=\frac{ke^{2}}{r}
\]
角动量量子化：\(mvr=n\hbar\)，即 \(v=n\hbar/(mr)\)。代入上式：
\[
m\cdot\frac{n^{2}\hbar^{2}}{m^{2}r^{2}}=\frac{ke^{2}}{r}\quad\Longrightarrow\quad r_{n}=\frac{\hbar^{2}}{mke^{2}}\,n^{2}
\]
括号里那堆常量就是玻尔半径 \(a_{0}=\hbar^{2}/(mke^{2})\approx5.29\times10^{-11}\) 米。再把 \(r_{n}\) 代回总能量 \(E=-ke^{2}/(2r)\)：
\[
E_{n}=-\frac{mk^{2}e^{4}}{2\hbar^{2}}\cdot\frac{1}{n^{2}}
\]
把电子质量、电荷、库仑常量、\(\hbar\) 的实验值代进去，系数算出 \(2.18\times10^{-18}\) 焦耳 \(=13.6\) 电子伏特——一个从纯常量组合里“凭空”算出的数字，恰好等于电离氢原子所需的能量。物理学的漂亮时刻往往长这样：你提出的规则预言了一个你从未测量过的数，而大自然签了字。
\end{mathbridge}

\begin{challenge}
玻尔模型还有一个深藏不露的成就：预言里德伯常量本身。把 \(E_{n}\) 代入跃迁公式 \(hc/\lambda=E_{n}-E_{m}\)，整理得到 \(1/\lambda=(E_{1}/hc)(1/m^{2}-1/n^{2})\)，也就是说 \(R=E_{1}/(hc)\)。右边全是基本常量，左边是光谱学测了几十年的数——两边在万分之几的精度内吻合。一个“编出来的规则”算对了六位有效数字，这不可能是巧合：玻尔模型抓住了某种真的东西。

另一头，算算旧理论输得有多惨。经典电磁学的拉莫尔公式说，加速电荷的辐射功率正比于加速度的平方。把绕核电子的向心加速度代进去，再除以它的总能量，得到的“辐射完所有能量”的特征时间是 \(10^{-11}\) 秒量级。同样的公式搬到宏观电路和无线电天线上却屡试不爽——旧理论不是错了，是越界了。判断一个理论何时越界，和记住理论本身一样重要。
\end{challenge}

\begin{figure}[htbp]
\centering
\begin{tikzpicture}[scale=1.0]
  \draw[thick] (0,0) -- (4.6,0) node[right=2pt] {\scriptsize \(n=1\ \ -13.6\ \mathrm{eV}\)};
  \draw[thick] (0,2.1) -- (4.6,2.1) node[right=2pt] {\scriptsize \(n=2\ \ -3.4\ \mathrm{eV}\)};
  \draw[thick] (0,2.5) -- (4.6,2.5) node[right=2pt] {\scriptsize \(n=3\ \ -1.5\ \mathrm{eV}\)};
  \draw[thick] (0,2.65) -- (4.6,2.65) node[right=2pt] {\scriptsize \(n=4\)};
  \draw[dashed] (0,3.1) -- (4.6,3.1) node[right=2pt] {\scriptsize \(n=\infty\ \ 0\)（电离）};
  \draw[-{Stealth},very thick,red!70!black] (1.3,2.5) -- (1.3,2.1);
  \node[red!70!black] at (2.35,2.28) {\scriptsize H\(\alpha\)：\(656\ \mathrm{nm}\) 红光};
  \draw[-{Stealth},very thick,blue!70!black] (3.6,2.1) -- (3.6,0);
  \node[blue!70!black] at (4.35,1.05) {\scriptsize \(122\ \mathrm{nm}\) 紫外};
\end{tikzpicture}
\caption{氢原子的能级阶梯（示意，纵轴按能量比例）。向下跳发光，向上跳吸光；档距决定颜色，级数越高档距越密。}
\end{figure}

\step{模型的边界}

\begin{mainline}
玻尔模型是物理学史上最成功的“错误模型”之一。先清点它成立的地盘：对一切单电子体系——氢原子、电离一次的氦离子、电离两次的锂离子——它给出的能级数值好得惊人，电离能、谱线位置全盘命中；它发明的“能级”“跃迁”“量子数”三个概念，被后来的量子力学原封不动地继承下来，至今仍是化学家和天体物理学家的日常语言。

再看它倒塌的地盘，一共四处，一处比一处根本。第一，多电子原子：给氦原子算能量，玻尔方法错得离谱，两个电子之间的相互作用让整个“指定轨道”的思路无从下手。第二，谱线强度：实验上有的线亮、有的线暗、有的干脆不出现，玻尔模型只能给位置，给不了亮度账。第三，精细结构：用更精细的光谱仪看，每条“线”其实是挨得极近的一小簇线，玻尔的整数梯子解释不了梯子缝里还有台阶。第四，也是最深的：轨道图像本身被证伪。1925 年之后量子力学建成，人们才明白电子根本不沿任何确定路线运动——不确定关系禁止电子同时拥有确定的位置和确定的动量，这不是仪器不够精细，而是“轨迹”这个概念在原子尺度根本失效。所以本书郑重声明：本章图里的圆轨道只是历史插图，氢原子里的电子不是一颗沿圆圈奔跑的小球；定态的真实图像是电子波函数在核周围形成的驻波花样，那是第 52 章的内容。

还要点名两种典型的误用。一是把基态想象成电子“停在核旁边休息”：若电子静止且位置确定，不确定关系会被当场违反，而且这种图像连“为什么不掉进核里”都回答不了——答案恰恰是量子化的波没有更低的驻波模式可用。二是把玻尔模型当成“氢原子的最终理论”去背：它是脚手架，不是大楼。盖楼时脚手架必不可少，楼成之后要拆；但拆掉的只是轨道图像，能级、跃迁、\(h\nu=\Delta E\) 全都留在楼里，至今承重。

它的适用范围可以这样记：问“氢原子能发什么颜色的光、电离要多少能量”，玻尔模型又快又准；问“为什么这条线亮那条线暗”“氦原子什么样”“分子怎么形成”，请直接翻去第 52 章和第 53 章，那里才有真正的量子力学答案。
\end{mainline}

\step{回头解释开场现象}

\begin{mainline}
现在回到燃气灶旁那撮盐。黄光来自钠原子最著名的一对跃迁：波长 \(589.0\) 和 \(589.6\) 纳米的两条线，肉眼分不开，合起来就是那滴纯黄。食盐在火焰里分解出钠原子，火焰的能量把钠的电子推上一档，电子跳回来时就把这笔账兑成黄光。烟花师傅的配方同理：锶的梯子在红色区有档距，钡在绿色区，铜在蓝色区——化学配方表其实是一份份能级表。

氢放电管的亮线、太阳光谱的暗线，谜底也一起揭开了。太阳内部的高温稠密气体发出连续彩虹，彩虹穿过较冷的太阳大气时，氢原子按自己的档距把 \(656\) 纳米、\(486\) 纳米这些特定波长吃掉，留下暗线；暗线和氢放电管的亮线位置分毫不差，因为同一架梯子，往上跳和往下跳用的是同样的档距。沿着这个逻辑，人类在 1868 年从日珥的光谱里发现了一条陌生的黄线，地球上当时没有任何元素对得上它——新元素氦就这样先在太阳上被“看见”，直到 1895 年才在地球的矿石气体里被拿到手里。元素可以在被触摸之前先被光指认，这是光谱学送给科学的礼物。

先猜一猜里的第一条也可以结案了。烧热的固体发出连续彩虹、颜色大体由温度决定，是因为固体里亿万个原子挤在一起，彼此的梯子互相踩得变形，档距被抹成密密麻麻的“能带”，于是什么颜色都能发（第 55 章细讲）；稀薄气体里的原子互不打扰，梯子档距分明，所以只发自己的颜色。同一条物理，两种面相，分界线是原子之间挨得多近。

工程上，这把“指纹锁”早已成了日常工具。原子吸收光谱仪让白光穿过样品烧成的蒸气，看彩虹里缺了哪几条线，就能报出样品里有什么元素、有多少——血铅筛查、自来水重金属检测、奶粉污染物抽检，靠的都是它，灵敏度可达十亿分之一量级。天体物理学家更进一步：恒星光谱的线系告诉他们恒星的成分和温度，谱线的整体移动（多普勒效应）告诉他们恒星在靠近还是远离，甚至替他们找到了围绕其他恒星转动的行星。你抬头看见的每一颗恒星，都寄来了一张写满原子账单的化验单。
\end{mainline}

\step{脑洞实验与挑战题}

\begin{thought}[把原子放大到一座体育场]
把氢原子的玻尔半径 \(5.3\times10^{-11}\) 米放大到五十米——大约半个足球场——放大倍数约 \(10^{12}\)。按同样倍数，原子核（一个质子，直径约 \(1.7\times10^{-15}\) 米）放大到约两毫米：场地正中央的一粒小玻璃珠。这粒珠子占了整个“体育场原子”百分之九十九点九七以上的质量，其余空间交给电子的量子态。你此刻坐着的椅子、握着这本书的手，“实心”的感觉几乎全部来自电子之间的电磁排斥和量子效应，而不是物质真的填满了空间。

再开一个更狠的脑洞。电子绕核辐射能量会坠毁——可地球绕太阳转，严格说也在辐射（引力波），为什么太阳系没塌？代入公式你会发现，地球轨道的引力辐射功率只有约两百瓦，相当于两只灯泡，而轨道的束缚能量高达 \(10^{33}\) 焦耳量级，照这个速度漏光要 \(10^{23}\) 年量级，比宇宙年龄还长十三万亿倍。同一个“辐射坠毁”的逻辑，在电磁世界一百亿分之一秒就执行完毕，在引力世界却慢得等于不存在。差别不在逻辑，在辐射功率与束缚能量的比例——物理学里，问“多大、多快、多少”和问“有没有”一样重要。
\end{thought}

\begin{mainline}
最后留几道题。重推理，轻计算，答案的关键都在思路上。
\end{mainline}

\begin{puzzles}
\item 一根氢放电管发出几条分立的亮线，一个白炽灯泡发出连续彩虹。两者都是“原子被加热后发光”，为什么光谱形态完全不同？思路提示：先别问发光，先问“这些原子彼此挨得多近”；孤立原子的梯子档距分明，亿万把梯子挤在一起时会发生什么？把答案留到第 55 章验证。
\item 一个氢原子的电子处在 \(n=3\) 的能级上，它最多能发出几种不同频率的光？如果是一亿个都处在 \(n=3\) 的氢原子呢？思路提示：先列出所有可能的跳法（\(3\to1\)、\(3\to2\)、\(2\to1\)），再注意单个原子跳一次就结束了，级联跳法里第二次跳跃是从哪一级出发的。单个原子与一大群原子的“最多”并不一样。
\item 十九世纪的天文学家在星光里发现一套神秘谱线，一度以为找到了新元素，后来才知道那是电离氦（失去一个电子的氦离子）。用本章的梯子语言说明：为什么氦离子的谱线会和氢的谱线“穿插”出现，以至人们会上当？思路提示：氦离子也是单电子体系，但核电荷是氢的两倍；想想 \(E_{n}\) 公式里核电荷藏在哪几项里（库仑力里的 \(e^{2}\) 一项会变成什么），能级整体会变成氢的几倍，再想想“四倍”遇上“\(1/n^{2}\)”会发生什么巧合。
\item 玻尔说定态中的电子不辐射，麦克斯韦说加速的电荷必辐射——这对矛盾在 1913 年是靠“立规矩”压下去的。量子力学真正是怎么收场的？思路提示：问题的陷阱在“绕核跑的小球”这个图像；如果定态里真正不随时间变化的是电子的电荷分布（驻波），那么“加速”这个前提还存在吗？再想想基态：为什么不存在比 \(n=1\) 更低的驻波模式。
\end{puzzles}

\begin{mainline}
本章结尾，请把镜头拉远。三种原子模型——实心球、枣糕、核式加梯子——两种进了博物馆，但它们不是“错误”,而是各自精度范围内的真相：道尔顿模型在化学反应的精度内至今正确，枣糕模型逼出了卢瑟福的炮弹，玻尔的梯子至今还在帮化学家估算。科学的前进不是新理论羞辱旧理论，而是每一次碰撞都把“成立的边界”画得更清楚。下一章，我们给氢原子换上真正的量子引擎：薛定谔方程。你会看到那架梯子在波动方程里自动长出来，而“轨道”这个词将获得一个和行星毫无关系的新含义。
\end{mainline}
